% id: 0707
% book: RPM 공통수학2 · 07 명제
% section: 유형 15 충분조건, 필요조건과 명제의 참, 거짓
% figure: none
% answer: ④
% answer_source: 답지
% variant_level: 0 (원본 전사)
% 이 파일은 build.mjs 가 items.json 에서 만든다. 고칠 때는 items.json 을 고친다.
\smash{\raisebox{1.5mm}{\dmkichul{RPM 공통수학2 0707번 · 105쪽}}}
세 조건 $p$, $q$, $r$에 대하여 두 명제 $p\longrightarrow q$, $\sim r\longrightarrow \sim q$가 모두 참일 때, \textbf{보기}에서 항상 옳은 것만을 있는 대로 고른 것은? \bogi{ㄱ. $q$는 $p$이기 위한 필요충분조건이다.\\ ㄴ. $r$는 $q$이기 위한 필요조건이다.\\ ㄷ. $p$는 $r$이기 위한 충분조건이다.}
\begin{choices32}{ㄱ}{ㄴ}{ㄷ}{ㄴ, ㄷ}{ㄱ, ㄴ, ㄷ}\end{choices32}
